\documentclass[9pt]{extarticle}
\usepackage[margin=0.65in, headheight=14pt]{geometry}
\usepackage{enumitem}
\usepackage{forest}
\usepackage{hyperref}
\usepackage{xcolor}
\usepackage{fancyhdr}
\usepackage{titlesec}

% === Style Guide (for AI-assisted updates) ===
%
% Tree layout
%   - forest package, monospace font, dot-and-line connectors
%   - each nesting level indented 19.5pt
%   - 9pt body, 0.65in margins, fancy header/footer
%
% Annotations
%   - \cmt{...}  → grey (#6A737D) italic — directory descriptions
%   - \fcmt{...} → orange (#FF9966) italic — file descriptions
%
% Writing style
%   - one short phrase per item: capture the high-level responsibility
%   - directories: what the folder contains or who it serves
%   - files: what the file does, not how it works
%
% Expansion rules
%   - docs/ directories: never expand, just annotate the directory
%   - shared/configs/: never expand, just annotate the directory
%   - results/ subdirs: list subdirectory names with annotations, don't list individual result files
%   - everything else: expand and annotate each file

% Unless explicitly prompted, do not modifiy p3, p4 directory structure

\definecolor{dircmt}{HTML}{6A737D}
\definecolor{filecmt}{HTML}{FF9966}
\newcommand{\cmt}[1]{\normalfont{\textcolor{dircmt}{\textit{#1}}}}
\newcommand{\fcmt}[1]{\normalfont{\textcolor{filecmt}{\textit{#1}}}}
\hypersetup{colorlinks=true, linkcolor=blue, urlcolor=blue}
\setlength{\parindent}{0pt}
\setlength{\parskip}{3pt}
\titlespacing*{\section}{0pt}{8pt}{3pt}
\titlespacing*{\subsection}{1.5em}{6pt}{2pt}
\setlist[itemize]{leftmargin=1.5em}

\title{\vspace{-1.2cm}\huge\textbf{A3: Directory Structure}\vspace{-0.5em}}
\author{\normalsize Autobook}
\date{\vspace{-1.5em}}

\pagestyle{fancy}
\fancyhf{}
\fancyhead[L]{\small A3: Directory Structure}
\fancyhead[R]{\small Autobook}
\fancyfoot[C]{\small \thepage}
\renewcommand{\headrulewidth}{0.4pt}
\renewcommand{\footrulewidth}{0pt}

\begin{document}
\maketitle
\thispagestyle{fancy}

How the \texttt{a3/} folder is organized inside the AI-Accountant repo.

\subsection*{Layout}

\begin{forest}
  for tree={
    font=\ttfamily,
    grow'=0,
    child anchor=west,
    parent anchor=south,
    anchor=west,
    calign=first,
    edge path={
      \noexpand\path [draw, -{|}] (!u.south west) +(7.5pt,0) |- node[fill,inner sep=1.25pt] {} (.child anchor)\forestoption{edge label};
    },
    before typesetting nodes={
      if n=1 {insert before={[,phantom]}} {}
    },
    fit=band,
    before computing xy={l=19.5pt},
    s sep=4pt,
  }
  [a3/
    [docs/ \cmt{teammate onboarding and shared references}]
    [shared/ \cmt{pipeline infrastructure used by all parts}
      [configs/ \cmt{smoke-test YAMLs for pipeline validation}]
      [scripts/ \cmt{Modal pipeline code}
        [main.py \fcmt{local entrypoint: reads YAML config, launches Modal run}]
        [shared/
          [infra.py \fcmt{Modal app, image, volume, and secrets definitions}]
          [helpers.py \fcmt{pure utilities: git checkout, checkpoint discovery, eval parsing}]
        ]
        [runners/
          [run.py \fcmt{orchestrator: dispatches training stages and evals in parallel}]
          [train.py \fcmt{runs nanochat training on a Modal GPU}]
          [evaluate.py \fcmt{runs BPB, CORE, and custom evals on a Modal GPU}]
        ]
      ]
    ]
    [p2/ \cmt{picochat ablation}
      [configs/
        [p2\_baseline.yaml \fcmt{vanilla picochat baseline}]
        [p2\_swiglu.yaml \fcmt{SwiGLU activation ablation}]
        [p2\_layerscale.yaml \fcmt{LayerScale ablation}]
        [p2\_swiglu\_layerscale.yaml \fcmt{combined SwiGLU + LayerScale ablation}]
        [p2\_*\_eval\_bpb\_core.yaml \fcmt{eval-only configs for each ablation}]
        [p2\_swiglu\_layerscale\_smoke.yaml \fcmt{quick smoke test for combined variant}]
      ]
      [training-troubleshooting.md \fcmt{notes on training issues encountered}]
    ]
    [p3/ \cmt{continued on next page}]
    [p4/ \cmt{continued on next page}]
  ]
\end{forest}

\newpage

\subsection*{Layout (continued)}

\begin{forest}
  for tree={
    font=\ttfamily,
    grow'=0,
    child anchor=west,
    parent anchor=south,
    anchor=west,
    calign=first,
    edge path={
      \noexpand\path [draw, -{|}] (!u.south west) +(7.5pt,0) |- node[fill,inner sep=1.25pt] {} (.child anchor)\forestoption{edge label};
    },
    before typesetting nodes={
      if n=1 {insert before={[,phantom]}} {}
    },
    fit=band,
    before computing xy={l=19.5pt},
    s sep=4pt,
  }
  [a3/
    [p3/ \cmt{context window extension}
      [docs/ \cmt{literature review, training config, eval design, final spec}]
      [configs/
        [p3\_baseline.yaml \fcmt{three-stage context extension experiment}]
        [p3\_eval\_only.yaml \fcmt{re-run evals on existing checkpoints}]
        [p3\_core\_evals.yaml \fcmt{CORE-only eval variant}]
        [p3\_from\_p2\_template.yaml \fcmt{template to eval P2's best model on P3 evals}]
      ]
      [evals/
        [context\_length\_eval.py \fcmt{positional perplexity eval on PG19 test split}]
      ]
      [results/ \cmt{run outputs}
        [core\_csvs/ \cmt{CORE per-task accuracy CSVs}]
        [custom\_evals/ \cmt{positional perplexity JSONs}]
        [figures/ \cmt{plots and summary tables}]
        [tex/ \cmt{compiled writeup PDF}]
      ]
      [analysis/
        [generate\_plots.py \fcmt{builds figures and tables from fetched results}]
      ]
    ]
    [p4/ \cmt{scaling laws and emergent abilities}
      [configs/
        [p4\_nano\_baseline\_fp8.yaml \fcmt{nanochat baseline with FP8 training}]
        [p4\_nano\_swiglu\_fp8.yaml \fcmt{nanochat SwiGLU variant with FP8 training}]
        [p4\_pipeline\_smoke\_20min.yaml \fcmt{quick smoke test for nano-scale runs}]
        [p4\_*\_eval\_only.yaml \fcmt{eval-only configs for each nano model}]
        [p4\_template.yaml \fcmt{template YAML for scaling runs}]
      ]
      [scripts/
        [compare\_emergent\_abilities.py \fcmt{finds questions nanochat answers but picochat cannot}]
      ]
      [results/ \cmt{run outputs}
        [final/ \cmt{final eval CSVs and JSONs per model}]
      ]
      [emergent\_abilities\_results.txt \fcmt{emergent abilities comparison output}]
    ]
  ]
\end{forest}

\end{document}
